\label{cap:conclusao}

Este trabalho investigou o estado da arte sobre a operacionalização de princípios éticos em sistemas de Inteligência Artificial no contexto da Engenharia de Software, com foco em mecanismos práticos capazes de apoiar desenvolvedores na verificação da conformidade ética ao longo do ciclo de vida do software. Para isso, foi conduzido um Mapeamento Sistemático da Literatura (MSL), abrangendo o período de 2021 a 2025, conforme as diretrizes metodológicas de Petersen et al.~\cite{Petersen2015}. Esse mapeamento constituiu a primeira fase da abordagem metodológica adotada, fundamentando-se na necessidade de consolidar os princípios éticos mais prevalentes na literatura e de identificar lacunas técnicas que orientassem o desenvolvimento dos artefatos propostos nesta dissertação.

Os resultados do MSL mostraram que a literatura sobre ética aplicada à IA concentra-se principalmente em três princípios centrais: \textit{Fairness}, \textit{Accountability} e \textit{Transparency}, identificados em 81,4\%, 55,9\% e 59,3\% dos estudos analisados, respectivamente. Esses princípios não apenas estruturam o núcleo conceitual das investigações recentes, como também são aqueles para os quais há maior maturidade em termos de métodos, \textit{frameworks} e ferramentas de apoio. Ainda assim, a análise detalhada revelou que a operacionalização ética permanece fragmentada, com predominância de abordagens voltadas para as fases de \textit{Requirements}, \textit{Design} e \textit{Implementation}, enquanto a fase de \textit{Testing} apresenta baixa cobertura (27,1\%), mesmo sendo crítica para a detecção de falhas éticas antes da implantação.

De modo geral, os achados do mapeamento evidenciam que, embora exista um corpo crescente de métodos, guias e ferramentas, ainda se carece de soluções integradas e automatizadas que permitam verificar, de forma sistemática e escalável, violações éticas em modelos de IA. Entre as principais lacunas identificadas, destacam-se: (i) ausência de mecanismos de rastreabilidade que conectem decisões éticas ao longo do ciclo de vida do software; (ii) escassez de ferramentas capazes de operar em ecossistemas reais de desenvolvimento (IDEs, pipelines de CI/CD, plataformas de MLOps); (iii) sub-representação de abordagens aplicadas à fase de testes; e (iv) falta de evidências empíricas robustas sobre a adoção industrial de práticas éticas. Tais lacunas, detalhadas no Capítulo~\ref{cap:rsl} e contextualizadas como motivação no Capítulo~\ref{introducao}, justificam a proposta desta dissertação e orientam a construção das \textit{fuzzing engines} dedicadas aos princípios de \textit{Transparency}, \textit{Fairness} e \textit{Accountability}.

Com base no diagnóstico fornecido pelo MSL, a pesquisa avançou para sua segunda fase, dedicada à identificação e seleção de riscos éticos passíveis de serem operacionalizados por meio de técnicas de \textit{fuzzing}, conforme detalhado na Seção~\ref{sec:criterios_selecao}. Essa etapa permitiu estabelecer uma taxonomia de riscos testáveis que fundamenta a especificação das engines: dos 11 riscos identificados, 6 demonstraram potencial pleno de implementação via \textit{fuzzing} automatizado --- evidenciado pela especificação técnica e pela avaliação empírica apresentadas no Capítulo~\ref{cap:proposta} ---, enquanto 5 apresentam limitações intrínsecas decorrentes de sua natureza organizacional, normativa ou dependência de julgamento humano especializado, requerendo métodos complementares de avaliação como auditorias qualitativas, estudos com usuários e análise por especialistas em ciências sociais. A especificação técnica adota um formato padronizado próprio, definindo para cada risco implementável: técnica de \textit{fuzzing} adequada, oráculo formal, critérios de aprovação/reprovação e métricas de avaliação.

A análise integrada de métodos de auditoria ética, ferramentas de \textit{fairness}, recursos de explicabilidade e abordagens de governança, detalhada na Seção~\ref{sec:metodos-correlatos}, revelou oportunidades concretas para a aplicação de \textit{fuzzing} como técnica inovadora para o teste de princípios éticos em IA, especialmente por sua capacidade de explorar automaticamente entradas diversificadas, detectar comportamentos inesperados e estimular condições extremas que podem revelar falhas éticas ocultas.

Portanto, as conclusões deste trabalho indicam que há uma lacuna clara e relevante na literatura e na prática da Engenharia de Software quanto à disponibilidade de ferramentas automatizadas para a avaliação sistemática de princípios éticos em IA. A partir desse diagnóstico, as etapas subsequentes da pesquisa concentraram-se em: (i) implementar \textit{fuzzing engines} dedicadas aos princípios de \textit{Transparency}, \textit{Fairness} e \textit{Accountability}; (ii) integrar as engines à arquitetura modular do \textit{framework}; e (iii) conduzir uma avaliação empírica com os modelos GPT (OpenAI), Gemini (Google) e DeepSeek, a fim de verificar a eficácia das engines e gerar evidências sobre a prevalência de falhas éticas em sistemas de IA contemporâneos. A avaliação empírica apresentada no Capítulo~\ref{cap:proposta}, conduzida sobre 8.880 pares avaliados e aproximadamente 14.160 chamadas de API, revelou perfis éticos distintos entre os provedores: taxas elevadas de falha em \textit{Fairness} (R-F1: 79,2\% global; R-F2: 55,4\% global; R-F4: 12,3\% global) e na detecção de vieses ocultos (R-T2: 70,0\% global, intervalo inter-provedor 65--75\%), com resultados moderados em \textit{Accountability} (R-A2: 30,8\% global, intervalo inter-provedor 27--33\%) e em \textit{Transparency} explicativa (R-T1: 12,8\% global, com forte heterogeneidade entre cenários: 22--27\% em explicação direta versus 0--1,5\% em meta-explicação). Agrupando por princípio, \textit{Transparency} apresentou a maior taxa global de falha (39,9\%), seguida de \textit{Accountability} (30,8\%) e \textit{Fairness} (30,1\%).

Os resultados obtidos incluíram a disponibilização de ferramentas modulares e reutilizáveis, integráveis à arquitetura do \textit{framework}, um conjunto expressivo de casos de teste para princípios éticos, métricas de conformidade que permitem análises em múltiplos níveis de granularidade --- da taxa de detecção individual por engine até o \textit{failure rate} consolidado por risco, modelo e princípio --- e diretrizes para adoção prática. O \textit{framework} foi disponibilizado como artefato público, acompanhado de uma interface gráfica que permite a configuração e execução de campanhas de teste de forma interativa. Esses artefatos contribuem tanto para avanços científicos quanto para a melhoria dos processos de desenvolvimento e auditoria ética de sistemas de IA na indústria.